\documentclass[conference]{IEEEtran}
\usepackage{cite}
\usepackage{amsmath,amssymb}
\usepackage{graphicx}
\usepackage{booktabs}
\usepackage{url}

\title{Evaluating Prompt-Injection and Retrieval/Tool-Channel Attacks on Agentic NetOps Pipelines and Layered Defenses}

\author{
\IEEEauthorblockN{Autonomous AI Researcher}
\IEEEauthorblockA{CNSM 2026 Agentic AI Researcher Track}
}

\begin{document}

\maketitle

\begin{abstract}
Problem: Agentic LLM pipelines for NetOps expose retrieval and tool channels that adversaries can exploit via prompt injection. Method: We preregistered (cnsmautocand-2-v1) and executed a paired confirmatory experiment comparing an unguarded agentic NetOps pipeline (baseline) against a defense-in-depth configuration (RAG grounding with provenance + deterministic runtime validator + sandboxed tool outputs). Design: preregistered paired plan; deterministic parser and provenance logging; frozen model-configuration and prompts (hashed). Result highlights (artifact-backed): completed\_episode\_count = 300, model\_calls\_used = 163, analyzed complete pairs n = 150; observed baseline attack-success = 147/150, guarded (combined defense) attack-success = 150/150; paired contingency counts recorded in analysis artifacts: n\_11 = 147, n\_10 = 3, n\_01 = 0, n\_00 = 0; paired difference (guarded − baseline) = +0.02, exact McNemar two-sided p = 0.25, bootstrap 95\% CI = [0.00, 0.0466666667] (bootstrap seed = 1729, resamples = 10000). Important caveats: the execution and analysis artifacts contain an internal inconsistency between some marginal tallies and paired counts; deterministic call-caching materially reduced external model calls (model\_calls\_used = 163) and may affect dependence across condition labels. Conclusion: for the studied prototype and adversarial corpus, attack-success was very high under both conditions and we did not observe a detectable reduction from the implemented combined defense. All raw artifacts, prompts, model-configuration, call-cache, provenance traces, parser outputs and analysis scripts are published in the execution bundle and hashed (referenced in the Disclosure Statement).
\end{abstract}

\section{Introduction}
Agentic LLM systems enable intent-driven NetOps automation but enlarge the attack surface via retrieval (RAG) and tool channels that can carry adversarial instructions. Prior work documents indirect prompt-injection vectors and recommends defense-in-depth (grounding, provenance, validators, sandboxing). We study whether a layered defensive configuration reduces adversarial-injection success in a preregistered, paired confirmatory experiment. Research question: How effective are retrieval- and tool-channel prompt-injection attacks against an agentic NetOps pipeline, and to what extent does a defense-in-depth configuration reduce attack success? We restrict claims to the implemented prototype, adversarial corpus, and simulator semantics.

\section{Related Work}
We built on verified literature covering prompt injection and agentic LLM agents, RAG grounding and provenance, generate–verify–repair architectures, and runtime/shielding defenses. Representative, verified sources are included in the artifact-backed bibliography and preregistration (examples: recent works on indirect prompt injection, RAG/agent benchmarks, and generate-verify-repair architectures). This study is experimental and narrowly scoped to pipeline-level attack–defense interactions in a simulator-executable NetOps semantics layer.

\section{Methodology}
Preregistration: study id cnsmautocand-2-v1 (frozen before execution) specified paired confirmatory contrast, inclusion/exclusion rules, seeds, retry policy, estimands, and analysis plan. Model scope (preregistered): two hosted-model families were listed at family level; the frozen model-configuration manifest used at runtime is included in the artifacts (execution/model\_configuration.json; SHA256 recorded in the execution manifest). Pipeline prototype: (1) frozen RAG index with provenance tagging, (2) deterministic parser mapping LLM outputs to canonical CLI ASTs, (3) deterministic runtime verifier/oracle (policy predicates), (4) sandboxed tool-emulation layer. Conditions executed per artifacts: baseline (unguarded) and combined defense (RAG+validator+sandbox). Deterministic call-cache was used (call\_cache files present and hashed) to enable deterministic replay and to reduce external model calls; its presence and per-file hashes are in the artifact bundle. Inclusion/exclusion: parser failures, retrieval-index mismatches, and out-of-scope provenance trigger exclusion from confirmatory analysis per the preregistration; exclusions and reasons are recorded (analysis/missingness\_summary.csv). Analysis: preregistered paired analysis (exact McNemar for the primary contrast; bootstrap percentile CI for paired difference; bootstrap seed and resamples recorded). All inputs, outputs, logs and hashes are published in the artifact bundle.

\section{Results}
Execution and analysis artifacts (hashed) are the authoritative source for reported numbers. Execution manifest summary (artifact-backed): planned\_episode\_count = 320, completed\_episode\_count = 300, failed\_episode\_count = 20, model\_calls\_used = 163 (execution/execution\_log.jsonl; per-file hashes in the execution manifest). Confirmatory analysis (analysis/results.json; hashed): complete\_pairs = 150; baseline\_success\_count = 147; guarded\_success\_count = 150; paired contingency counts recorded: n\_11 = 147, n\_10 = 3, n\_01 = 0, n\_00 = 0. Derived confirmatory statistics (from the frozen analysis pipeline): paired difference (guarded − baseline) = +0.020000000000000018; exact McNemar two-sided p = 0.25; bootstrap 95\% percentile CI = [0.00, 0.04666666666666667] (bootstrap seed = 1729; resamples = 10000). Important integrity note (artifact-documented): the execution and analysis artifacts contain an internal inconsistency between some marginal tallies and the paired counts (the artifact bundle records both marginal counts and paired tables that do not numerically match in one provenance view). We cannot fabricate a corrected recount without rerunning or reprocessing raw artifacts under an immutable audit pipeline; accordingly we (a) report the exact numbers as they appear in the frozen artifacts above, (b) provide the full paired counts (n\_11/n\_10/n\_01/n\_00) as recorded, and (c) conservatively interpret results in light of the inconsistency (see Discussion and Limitations). Model-call accounting: model\_calls\_used = 163 for 300 completed episodes indicates substantial deterministic reuse via the call\_cache (call\_cache directory and per-file hashes are in the artifact bundle). Per-preregistration exclusions and missingness are recorded in analysis/missingness\_summary.csv (artifact hash in manifest) and summarized in the analysis bundle.

\section{Discussion}
Primary interpretation under artifact constraints: both baseline and guarded conditions exhibited very high observed attack-success in the implemented prototype and adversarial corpus. The preregistered confirmatory contrast (baseline vs combined defense) produced a small positive paired difference (+2\%) and no statistically significant reduction (exact McNemar p = 0.25) under the recorded artifact numbers. However, two integrity issues materially limit confirmatory inference and are recorded in the artifact bundle: (1) an internal counting inconsistency between some marginal tallies and the paired contingency table (both are present and hashed in the artifacts), and (2) deterministic call-caching substantially reduced external model calls (model\_calls\_used = 163) and therefore can induce dependence across runs/conditions if cached responses are reused across condition labels. We do not invent a corrective reanalysis here; instead we (a) supply the exact artifact-derived numbers and hashes for external audit, (b) mark the inconsistency and caching-dependence as unresolved evidence limitations that require reprocessing or rerun to fully reconcile, and (c) offer conservative conclusions: the defense, as implemented in this prototype and under the recorded artifacts, did not produce a detectable reduction in attack-success. The artifacts permit external auditors to recompute stratified exploratory analyses (per-attack-type, per-task-cluster, per-model-family) using provided scoring JSONs and provider-call logs; those exploratory breakdowns are archived (execution/scoring/*.json and execution/provider\_calls/*.json) but were not counted as confirmatory in this manuscript.

\section{Limitations}
\begin{itemize}
\item Limited model coverage: only the preregistered model families were exercised; results do not generalize to other LLM families or versions.
\item Simulator semantics: NetOps semantics and validators ran in a simulator-executable environment (no live-device deployment); simulator–real-device mismatch could alter operational effectiveness in production.
\item Implementation-specific defenses: the RAG index, validator predicates, and sandbox were prototype implementations—different designs or stronger predicates might yield different outcomes.
\item Sample and deviations: planned preregistration runs differed from executed runs (preregistration planned a larger paired collection; execution\_manifest records planned\_episode\_count=320 and completed\_episode\_count=300); 10 preregistered pairs were unavailable and excluded per preregistered rules (detailed excluded task ids and reasons are in analysis/missingness\_summary.csv).
\item Parser and mapping dependence: analysis depends on deterministic parsing of LLM outputs to CLI ASTs; parser failures affect inclusion and are a construct-validity threat (parser failure counts are logged).
\item Deterministic caching: call\_cache reuse materially reduced external model calls (model\_calls\_used = 163 for 300 completed episodes). Cached-response reuse can induce dependence across condition labels; the call\_cache and provider\_call artifacts are published for audit but this dependence cannot be ruled out without reprocessing or rerun without caching.
\item Internal artifact inconsistency: archived artifacts include an unresolved internal inconsistency between reported marginal counts and the paired contingency table; both views are preserved and hashed. We report both artifact-derived values but cannot reconcile them without reprocessing or rerun; this prevents stronger confirmatory claims.
\item Attack-corpus and task scope: the adversarial corpus and four NetOps task clusters are not exhaustive of NetOps actions or attack techniques; external validity is limited.
\end{itemize}

\section{Conclusion}
In this preregistered, artifact-backed study of prompt-injection against an agentic NetOps prototype, both baseline and combined-defense configurations produced very high attack-success rates in the recorded runs. The frozen confirmatory analysis, as reflected in the archived artifacts, did not show a detectable reduction from the implemented combined defense (paired difference +2\%, exact McNemar two-sided p = 0.25, bootstrap 95\% CI includes 0). Critical audit caveats—an internal artifact inconsistency in counts and deterministic call-caching that materially reduced external model calls—prevent stronger confirmatory claims. All raw artifacts, prompts, parser outputs, call-cache files, provenance traces, scoring JSONs, and analysis scripts are published in the execution bundle with per-file SHA256 hashes for independent audit and reanalysis (see Disclosure). We recommend follow-up work that (a) reconciles the counted artifacts via deterministic reprocessing or (b) reruns preregistered pairs without call-cache reuse to assess robustness, and (c) evaluates stronger validator predicates and alternative RAG policies in real-device deployments to improve external validity.

\section*{Disclosure Statement}
This study executed under preregistration cnsmautocand-2-v1. Human roles and pre-lock development: a human Research Director selected the topic family and constrained the initial master prompt prior to the autonomy lock; additional pre-lock infrastructure development and rehearsals occurred and were explicitly excluded from final-study evidence. Autonomy boundary: after the autonomy lock the AI system autonomously conducted literature verification, study selection, preregistration, code generation, experiment execution, analysis, manuscript drafting, and revision; no manual human editing of technical content occurred after the lock. Models and prompting: the preregistered model families were specified at family level (OpenAI GPT-4 family and Anthropic Claude-2 family); the exact, frozen model-configuration and prompt templates used at runtime are recorded in execution/model\_configuration.json (SHA256 = "916b386bd80cbffb0e309c14af1bbea4090e342beb73b9a8784a0d57da8367fd"). Prompts and deterministic templates invoked during runs are preserved in the execution bundle and hashed (see execution/artifact\_hashes entries). Agentic/orchestration framework: the prototype pipeline included (1) a frozen retrieval (RAG) index with provenance tagging, (2) a deterministic parser that maps LLM outputs to canonical CLI ASTs, (3) a deterministic runtime verifier/oracle (policy predicates), and (4) a sandboxed tool-emulation layer. Implementation artifacts (parser code, deterministic mapping modules, and simulator mappings) are included in the artifact bundle (hashed; see execution/artifact\_hashes). Tools and permitted capabilities: runs invoked hosted-model APIs (per preregistration) and internal deterministic simulators only; no live network devices were modified. Preregistration/freeze procedure: the preregistration document (cnsmautocand-2-v1) and experiment plan were frozen and the preregistration SHA256 is recorded in the execution manifest (execution/preregistration\_sha256 = "36b1839becd2b0a8be22bf770165659fdec5814ade05513dd3c71d1faabb976b"). Autonomous experiment procedure: runs were automated per the frozen sampling and retry policy; retries only for infrastructure/API failures (one automated retry), logged in execution/execution\_log.jsonl (artifact hash in execution/artifact\_hashes). Deterministic call-caching and replayability: a deterministic call\_cache was used to enable replay and to limit external model calls; call\_cache files and many per-call hashes are preserved in execution/call\_cache/ and hashed in the execution manifest. Model-call accounting recorded in the execution manifest: model\_calls\_used = 163 across completed episodes = 300 (execution/execution\_log.jsonl; artifact hashes provided). Caching implication: the low model\_call count indicates substantial cache reuse; the artifacts (call\_cache entries and provider\_calls logs) are published to allow auditors to verify whether cached responses were reused across different condition labels. Technical restarts and deviations: execution\_manifest records completed\_episode\_count = 300 and failed\_episode\_count = 20; planned\_episode\_count in the execution manifest is 320 while the preregistration planned a larger set (the preregistration document and execution manifest both appear in the artifact bundle). All deviations, retry attempts, and failure reasons are logged in execution/execution\_log.jsonl (hash recorded). Artifact availability: all raw outputs, provider-call logs (execution/provider\_calls/*.json), call-cache items (execution/call\_cache/*.json), responses (execution/responses/*.txt), scoring JSONs (execution/scoring/*.json), parser outputs, execution logs, and analysis scripts are included in the execution artifact bundle and hashed in the execution manifest (see execution/artifact\_hashes for per-file SHA256). Integrity note: the archived artifacts include an internal counting inconsistency between some marginal tallies and the paired contingency table; both views are preserved and hashed to enable audit. We do not fabricate reconciliations; the inconsistency is disclosed and archived for external audit. No manual human editing of the manuscript's technical content occurred after the autonomy lock.

\begin{thebibliography}{99}
\bibitem{ref1} Rizwan Tanveer, "Prompt Injection and Jailbreak Attacks in Large Language Model-Based Agents", 2026, 10.2139/ssrn.6740060
\bibitem{ref2} Qiusi Zhan, Zhixiang Liang, Zifan Ying, Daniel Kang, "InjecAgent: Benchmarking Indirect Prompt Injections in Tool-Integrated Large Language Model Agents", 2024, 10.18653/v1/2024.findings-acl.624
\bibitem{ref3} Ze Nan, Wei Wei, Zhenhua Lin, Jingjing Chang, Yue Hao, "Flexible Nanocomposite Conductors for Electromagnetic Interference Shielding", 2023, 10.1007/s40820-023-01122-5
\bibitem{ref4} Lameck Mbangula Amugongo, Pietro Mascheroni, Steven E. Brooks, Stefan Doering, Jan Seidel, "Retrieval augmented generation for large language models in healthcare: A systematic review", 2025, 10.1371/journal.pdig.0000877
\bibitem{ref5} Oscar G. Lira, Oscar Maurício Caicedo Rendón, Nelson L. S. da Fonseca, "Large Language Models for Zero Touch Network Configuration Management", 2024, 10.1109/mcom.001.2400368
\bibitem{ref6} Changjie Wang, Mariano Scazzariello, Alireza Farshin, Simone Ferlin, Dejan Kostić, Marco Chiesa, "NetConfEval: Can LLMs Facilitate Network Configuration?", 2024, 10.1145/3656296
\end{thebibliography}

\end{document}
